\section*{Discussion}
% I'm reclaiming em dashes for humans. I'll disclose AI use as honestly as I can but I typed those three dashes with my own two hands
Here, we show that direct activation of ER2 neurons, co-activation of ER2/ER4 neurons, and inhibition of ER4d neurons are each sufficient to promote longevity in adult flies of both sexes.
% deleted "under standard laboratory conditions"
These effects are elicited by manipulating very small neuronal populations---in some cases, fewer than ten neurons per hemisphere \citep{omoto2018,hulse2021}---underscoring the sensitivity of organismal aging to defined circuit states.
Our findings establish that the EB can act instructively to modulate aging rather than serving only as a permissive gate for lifespan-altering sensory information.
Previous work showed that ER2 neurons are required for death perception to accelerate aging \citep{gendron2023} and, notably, that constitutive activation of ER2/ER4 neurons using heat-gated TRPA1 modestly reduced lifespan \citep{gendron2023}.
Read together with the data reported here, these contrasting outcomes suggest that the timing and persistence of EB activity help define ER-neuron-dependent internal states and their effects on aging, consistent with the context-sensitivity of valence that we observed.

We also find that manipulations of different EB populations likely modulate aging through distinct molecular mechanisms.
Bulk head transcriptomics after one week of treatment revealed that ER2 activation, ER4d inhibition, and ER2/ER4 co-activation produced largely non-overlapping transcriptional programs.
The differences were not merely quantitative: where the manipulations converged on shared pathways, they often did so in opposite directions, with starvation- and catabolism-related genes upregulated by ER2 activation but downregulated by ER2/ER4 co-activation.
Rather than resolving immediately to a common state, then, distinct activity patterns across ER2 and ER4 populations appear to be proximally interpreted as distinct internal states that each extend lifespan. 

Our experiments also revealed, perhaps fortuitously, that ER4d inhibition may enhance the ability to buffer moderate stress.
Green light alone shortened lifespan and broadly induced heat-shock proteins (e.g., \textit{Hsp23}, \textit{Hsp26}, \textit{Hsp27}, \textit{Hsp70Bb}) in both control flies and flies expressing GtACR1.
ER4d silencing produced few detectable genotype-specific changes apart from attenuating this stress signature.
Because temperatures in our ventilated rigs were not meaningfully elevated above those of dark-aged controls, heat is unlikely to explain these effects.
Fly vision is more sensitive to green than to red light \citep{salcedo1999}, and certain wavelengths can exert physical effects on the brain \citep{song2022} and shorten lifespan \citep{johnson2023}; the intensity and duration of exposure required for deep brain penetration in our experiments \citep{wolff2015} likely rendered green light a meaningful stressor.
Whether the attenuated molecular response in ER4d-silenced flies actively contributed to their longer lifespan, or instead reflects a reduced stress burden, remains to be determined, as does whether a comparable mechanism operates under more benign conditions.

The transcriptional data also point to specific pathways that may link EB activity to aging-regulatory physiology.
Following ER2 activation, the induction of \textit{foxo}, \textit{Thor}, \textit{InR}, \textit{dilp6}, and \textit{Cbs} suggests a coordinated engagement of nutrient-sensing, proteostatic, and sulfur-amino-acid metabolic processes previously implicated in lifespan control.
This state does not, however, resemble a canonical starvation program: nutrient-responsive IIS components were induced alongside ribosome biogenesis and translational programs that are not ordinarily co-regulated in this way \citep{essers2016}, suggesting that direct EB activation may uncouple pathways normally coordinated by nutrient availability.
Because prior work showed that ER2/ER4 neurons regulate \textit{dilp3} and \textit{dilp5} expression in insulin-producing cells during death perception \citep{gendron2023}, one plausible hypothesis is that EB-dependent internal states are translated into systemic aging effects through downstream neuroendocrine centers, though the present data do not establish this mechanism directly, and defining the relevant effector pathways will require future work.

The internal states established by EB manipulation were molecularly specific yet were not accompanied by the behavioral or physiological changes those states would ordinarily predict.
GO analysis revealed opposite transcriptional indicators of nutrient availability in the head: starvation-response processes were upregulated by ER2 activation (\figref{fig:seqfig}B) but downregulated by ER2/ER4 activation, which instead enriched for biosynthetic and storage-associated processes including amino acid synthesis and glycogen metabolism (\figref{fig:seqfig2}C), suggestive of a nutrient-replete state.
Neither signature was matched by a change in feeding.
The ER2 signature could have reflected reduced intake or a hunger-like state that should itself promote feeding, yet we observed neither; the nutrient-replete-like ER2/ER4 signature was likewise unaccompanied by any feeding change.
These transcriptional states are therefore unlikely to arise secondarily from altered consumption, and they argue against dietary restriction as the basis of lifespan extension.
The same dissociation extended to other coarse phenotypes: despite unchanged feeding, TAG stores were reduced in ER2 males and in both sexes after ER2/ER4 activation, yet these changes tracked neither lifespan extension nor starvation resistance, and negative geotaxis revealed no early correlate of the one-week transcriptional states.
Together, these results indicate that EB-dependent regulation of aging is not readily explained by gross changes in feeding, lipid storage, or locomotor performance.

The absence of feeding and physiological correlates did not, however, mean that EB manipulation was motivationally silent; each manipulation instead imposed a distinct, dissociable motivational state.
In the FLIC assay, ER2 activation in females was aversive, promoting rapid disengagement from the sucrose well without suppressing re-approach, whereas ER2/ER4 co-activation was positively valenced, increasing both the number of interactions and the frequency of return without prolonging individual feeding bouts.
In a feeding-independent place-preference assay, the positive valence of ER2/ER4 co-activation generalized across sex and stimulation regimes. By contrast, aversive states were more context-dependent: the negative valence of ER2 activation was restricted to females but held across stimulation regimes, while that of ER4d activation generalized across sex but only emerged under different stimulation regimes.
Together, these findings indicate that EB ring neurons do not contribute fixed, additive aversive or appetitive signals; rather, the motivational meaning of EB activity depends on population combination, assay context, sex, and stimulation conditions.

A particularly informative consequence of this combinatorial structure is that manipulations of opposite valence both extend lifespan: ER2 activation, which was aversive in females, and ER2/ER4 co-activation, which was appetitive, each reliably increased survival in males and females.
This is difficult to reconcile with models in which longevity is mediated by a particular motivational valence, or by any single downstream behavioral or physiological state that valence would normally engage.
It is also difficult to reconcile with a strict reading of the incoherent feedforward circuit in which ER2, ER4m, and ER4d neurons are thought to participate \citep{ishimoto2020}, under which ER2 activation, ER2/ER4 co-activation, and ER4d inhibition should converge on a common downstream state. This model is at odds with the divergence we observed in transcriptional profile, valence, and downstream physiology, although we cannot exclude convergence on a shared terminal mechanism farther downstream of the EB.
A more parsimonious interpretation is that EB circuits influence aging through population-level configurations whose downstream consequences depend on the particular combination of neurons engaged, consistent with the EB's established role as an integrative rather than relay structure.
In this view, lifespan extension is not the signature of a single privileged EB activity pattern but a property of multiple distinguishable circuit states, each recruiting partially distinct downstream programs that nonetheless converge on slower aging, with the transcriptional divergence we observed implying distinct molecular routes to a shared outcome.

Molecular profiling of peripheral tissues revealed that ER2 activation drove coordinated systemic remodeling beyond the brain.
In bodies, ER2 activation suppressed proliferative and DNA repair programs while upregulating a vacuolar H$^{+}$-ATPase-dominated electrochemical homeostasis program not observed in heads, indicating tissue-specific rather than uniform physiological responses.
That this V-ATPase program was not accompanied by autophagy or trafficking pathways suggests a decoupling of organellar acidification from degradative and secretory functions.
The accompanying upregulation of inwardly rectifying potassium channels may help maintain cell membrane potential during increased organellar proton flux, while the enrichment of fatty acid and ether phospholipid synthesis genes suggests a remodeling of cellular and organellar membranes to support increased ion transport.

Metabolomic profiling reinforced this interpretation.
The depletion of free nucleotides and amino acids is consistent with the reduced proliferative burden evident in our transcriptional data, while reductions in choline, glycerophosphocholine, and acylcarnitines may reflect increased membrane remodeling activity.
Among upregulated metabolites, the accumulation of glycolytic and TCA-cycle intermediates, together with glutamine and arginine, is consistent with reduced consumption of these species by growth-related processes.
Serine and betaine serve as methyl donors, and their elevation, alongside depletions in SAH and choline and the upregulation of both \textit{Sam-S} and \textit{Gnmt} in head and body, indicates a restructuring of methionine cycle flux.
Although none of the core kynurenine pathway genes were upregulated in our RNA-seq dataset, the coordinated elevation of kynurenic acid, quinolinic acid, and anthranilic acid could account for the depleted tryptophan signature.
The convergence of transcriptomic and metabolomic data indicates that activation of the small population of ER2 neurons imposes an orchestrated, tissue-specific remodeling of peripheral physiology that is distinct from canonical starvation or dietary restriction.

Our finding that discrete or combined populations of ellipsoid body ring neurons can instructively modulate lifespan by generating distinct internal states carries potentially significant implications for understanding the neural basis of aging in humans.
The \textit{Drosophila} central complex, of which the ellipsoid body is a core component, shares strong neuroanatomical and functional homology with the vertebrate basal ganglia \citep{strausfeld2013}.
Both structures serve as integrative hubs that process sensory information, coordinate motor output, and critically contribute to the generation of motivational and affective states.
In particular, limbic basal ganglia output pathways, including ventral pallidal circuits, encode opposing appetitive and aversive motivational value and promote approach or avoidance behaviors according to internal state and motivational context \citep{stephenson-jones2020}.
Our findings raise the possibility that a similarly general principle applies to aging: circuit configurations that impose distinct motivational or physiological states may converge on long-term effects on organismal maintenance and survival.
More generally, these results suggest that lifespan may be modulated not only by what animals perceive, but by how neural state transforms the meaning of those perceptions into coordinated physiological output.
If analogous circuits in the human brain similarly encode valence states that influence systemic physiology, this would provide a mechanistic framework for understanding associations between psychological states, chronic stress, and aging in humans.